\documentclass[12pt]{article}
\usepackage[utf8]{inputenc}
\usepackage[english,russian]{babel}
\usepackage{amsmath}
\usepackage{amssymb}
\usepackage{graphicx}
\usepackage{geometry}
\usepackage{tocloft}
\renewcommand{\cftsecleader}{\cftdotfill{\cftdotsep}} % Точки для разделов

% Уменьшаем верхний отступ на титульной странице
\geometry{
    top=2cm,       % Верхнее поле
    bottom=2cm,    % Нижнее поле
    left=2cm,      % Левое поле
    right=2cm,     % Правое поле
}

\begin{document}

% ТИТУЛЬНЫЙ ЛИСТ
\thispagestyle{empty} % Убираем номер страницы и колонтитулы
\begin{center}
МИНИСТЕРСТВО НАУКИ И ВЫСШЕГО ОБРАЗОВАНИЯ

РОССИЙСКОЙ ФЕДЕРАЦИИ

\vspace{1cm}

Федеральное государственное автономное

образовательное учреждение высшего образования

«Национальный исследовательский Нижегородский государственный

университет им. Н.И. Лобачевского (ННГУ)»

\vspace{4cm}

\textbf{Лабораторная работа на тему:}

\vspace{0.5cm}

\textbf{«Стабилизация верхнего положения перевёрнутого
маятника»}

\vspace{3cm}

\begin{flushright}
\small
Выполнили: Малеев Кирилл, Сорока Сергей, \\
Муратов Егор, Жданкин Александр, \\
Осадчий Дмитрий, Кабанчиков Съебас, \\
Орехов Никита.\\
Группа: 3823Б1ПМоп3 \\
Подгруппа № 3
\end{flushright}

\vspace{7cm}

\textbf{\large Нижний Новгород}

\textbf{\large 2025}
\end{center}

% Теперь устанавливаем обычный стиль страниц для остальных страниц
\pagestyle{plain} % Обычные номера страниц

\newpage

% ОГЛАВЛЕНИЕ (СОДЕРЖАНИЕ)
\renewcommand{\contentsname}{СОДЕРЖАНИЕ:}

\tableofcontents

\newpage

\section{Введение}

 $\indent$Данная работа посвящена исследованию задачи стабилизации перевернутого маятника с подвижной точкой опоры. Задача удержания перевернутого маятника в верхнем положении является классической задачей теории управления. Несмотря на кажущуюся простоту, она служит удобной моделью
для отработки методов управления неустойчивыми объектами, так как сочетает в себе нелинейность
и высокую чувствительность к возмущениям.

Актуальность этой задачи обусловлена её широким применением в современной технике. Стабилизация перевернутого маятника является ключевым элементом в системах управления ракетаминосителями на старте, где необходимо удерживать вертикальное положение центра масс. Аналогичные принципы лежат в основе управления современными высокоманевренными истребителями с
передней горизонтальной плоскостью оперения (ПГО), а также в системах активной подвески автомобилей для гашения колебаний кузова. Более того, эта модель служит прототипом для балансирующих роботов (например, Segway) и является важным этапом в изучении мехатронных систем.

В данной работе рассматриваются две стратегии управления: первая обеспечивает стабилизацию
только угла отклонения маятника, вторая — полную стабилизацию (угол и положение точки опоры). В работе приведён вывод уравнений движения, анализ устойчивости с использованием критерия Рауса-Гурвица, а также результаты численного моделирования, подтверждающие теоретические
выводы.

\newpage

\section{Уравнение движения маятника}

\begin{figure}[h!]
\centering
\includegraphics[width=0.7\textwidth]{catjpg.jpg}
\caption{Кот(обыкновенный)}
\label{fig:risunok}
\end{figure}


Функция Лагранжа определяется как разность кинетической и потенциальной энергий:

\begin{gather}
L=K-U,
\end{gather}
где кинетическая энергия $K$ и потенциальная энергия $U$ заданы выражениями:
\begin{gather}
K = \frac{m}{2}[(\dot{u}^2-\dot{\varphi}l\cos(\varphi))^2+l^2\dot{\varphi}^2sin^2(\varphi)] = \frac{m}{2}(\dot{u}^2-2\dot{\varphi}\dot{u}\cos(\varphi)+l^2\dot{\varphi}^2),\nonumber\\
U=mglcos(\varphi).
\end{gather}
Полный лагранжиан:
\begin{gather}
L=\frac{m}{2}(\dot{u}^2-2\dot{\varphi}\dot{u}\cos(\varphi)+l^2\dot{\varphi}^2)-mglcos(\varphi).
\end{gather}
Уравнение Лагранжа:
\begin{gather}
\frac{d}{dt}\left(\frac{dL}{d\dot{\varphi}}\right)-\frac{dL}{d\varphi}=0.
\end{gather}
Найдем частные прлизводные:
\begin{gather}
\frac{dL}{d\varphi}= \frac{m}{2}2l\dot{u}\dot{\varphi}\sin(\varphi)+mgl\sin(\varphi).\nonumber\\
\frac{dL}{d\dot{\varphi}}=-\frac{m}{2}2l\dot{u}\cos(\varphi)+ml^2\dot{\varphi}.\nonumber\\
\frac{d}{dt}\left(\frac{dL}{d\dot{\varphi}}\right)=-ml\ddot{u}\cos(\varphi)+ml^2\ddot{\varphi}+ml\dot{u}\dot{\varphi}sin(\varphi)
\end{gather}
Подставим (5) в (4):
\begin{gather}
-ml\ddot{u}\cos(\varphi)+ml^2\ddot{\varphi}+ml\dot{u}\dot{\varphi}sin(\varphi)-ml\dot{u}\dot{\varphi}\sin(\varphi)-mgl\sin(\varphi)=0    \quad| :ml\nonumber\\
\ddot{u}\cos(\varphi)+l\ddot{\varphi}+\dot{u}\dot{\varphi}sin(\varphi)-\dot{u}\dot{\varphi}\sin(\varphi)-g\sin(\varphi)=0.\nonumber\\
\ddot{u}\cos(\varphi)+l\ddot{\varphi}-g\sin(\varphi)=0.\nonumber\\
\ddot{\varphi}-\frac{\ddot{u}}{l}\cos(\varphi)-\frac{g}{l}\sin(\varphi)=0.\nonumber\\
\boxed{\ddot{\varphi} - \frac{g}{l} \sin(\varphi) = \frac{\ddot{u}}{l} \cos(\varphi)} \qquad \text{(уравнение движения с управлением)}
\end{gather}



\newpage














\section{Формулки из п4 и п5}

Мне в падлу все в соло делать, так что я просто вам формул накатаю(не все), а вы уже сами как то справитесь.
\\\\
-------------------------------------

Для анализа устойчивости найдем решенеи уравнения в виде:
\begin{gather}
\varphi(t)=Ae^{\lambda t}\nonumber\\
\lambda^2Ae^{\lambda t}+\frac{b}{l}\lambda Ae^{\lambda t}+\frac{a-g}{l} Ae^{\lambda t}=0    \quad:Ae^{\lambda t}\nonumber\\
\lambda^2+\frac{b}{l}\lambda+\frac{a-g}{l}=0.\nonumber
\end{gather}

Введем обозначения:
\begin{gather}
\frac{b}{l}=2\delta, \quad \frac{a-g}{l}=\omega^2,\nonumber
\end{gather}
где $\delta$ - коэффицент затухания, а $\omega$ -собственная частота системы.
\begin{gather}
\lambda^2 +2\delta\lambda+\omega^2=0
\end{gather}

Корни характеристического уравнения находятся по формуле:
\begin{gather}
\lambda_{1,2} = \frac{-2\delta\pm\sqrt{4\delta^2-4\omega^2}}{2}.\nonumber\\
\lambda_{1,2} = -\delta \pm \sqrt{\delta^2-\omega^2}.
\end{gather}
-----------------------------------------------

Для устранения этого недостатка необходимо расширить закон управления, добавив обратную
связь по положению и скорости точки опоры. Расширенное управление имеет вид:
\begin{gather}
\ddot{u}=a\varphi+b\dot{\varphi} +cu+d\dot{u}.
\end{gather}
где $a,b,c,d$ - параметры регулятора.\\
Линеаризованное уравнение движения перевернутого маятница имеет вид:
\begin{gather}
\ddot{\varphi}- \frac{g}{l}\varphi=\frac{\ddot{u}}{l}.
\end{gather}
Подсталяя управление (9) в уравнение маятника (10), получаем:
\begin{gather}
\ddot{\varphi}-\frac{g}{l}\varphi=\frac{1}{l}(a\varphi+b\dot{\varphi}+cu+d\dot{u}).\nonumber\\
\ddot{\varphi}-\frac{g}{l}\varphi=\frac{a}{l}\varphi+\frac{b}{l}\dot{\varphi}+\frac{c}{l}u+\frac{d}{l}\dot{u}.
\end{gather}
Выражаем $\ddot{\varphi}$:
\begin{gather}
\ddot{\varphi}=\frac{a}{l}\varphi+\frac{b}{l}\dot{\varphi}+\frac{c}{l}u+\frac{d}{l}\dot{u}+\frac{g}{l}\varphi.\nonumber\\
\ddot{\varphi}=\frac{b}{l}\dot{\varphi}+\frac{a+g}{l}\varphi+\frac{c}{l}u+\frac{d}{l}\dot{u}.
\end{gather}

Таким образом, замкнутая система описывается двумя диффурунциальными уравнениями:

\begin{equation}
\begin{cases}
\ddot{\varphi} = \dfrac{b}{l} \dot{\varphi} + \dfrac{a+g}{l} \varphi + \dfrac{c}{l} u + \dfrac{d}{l} \dot{u} \\ 
\ddot{u} = a\varphi + b\dot{\varphi} + cu + du 
\end{cases}
\end{equation}

Для анализа устойчивости ищем решение в виде экспонент:

\begin{equation}
\varphi(t) = Ae^{\lambda t}, \quad u(t) = Be^{\lambda t}
\end{equation}

Подставляя (14) в систему (13), получаем:

\begin{equation}
\begin{cases}
\left(\lambda^2 - \dfrac{b}{l}\lambda - \dfrac{a+g}{l}\right)A - \left(\dfrac{c}{l} + \dfrac{d}{l}\lambda\right)B = 0 \\[1em]
\left(a + b\lambda\right)A + \left(-\lambda^2 + c + d\lambda\right)B = 0
\end{cases}
\end{equation}

Переносим все члены в левую часть:
\begin{equation}
\begin{cases} 
\left( \lambda^2 - \dfrac{b}{l} \lambda - \dfrac{a + g}{l} \right) A - \left( \dfrac{c}{l} + \dfrac{d}{l} \lambda \right) B = 0 \\ 
(a + b\lambda)A + (c + d\lambda - \lambda^2)B = 0 
\end{cases}
\end{equation}

Система (16) имеет ненулевое решение тогда и только тогда, когда определитель матрицы коэффициентов равен нулю:

\begin{gather}
\det 
\begin{pmatrix} 
\lambda^2 - \dfrac{b}{l} \lambda - \dfrac{a + g}{l} & -\left( \dfrac{c}{l} + \dfrac{d}{l} \lambda \right) \\ 
a + b\lambda & c + d\lambda - \lambda^2 
\end{pmatrix} = 0
\end{gather}

Вычисляем определитель:

\begin{gather}
\left( \lambda^2 - \frac{b}{l} \lambda - \frac{a + g}{l} \right) (c + d\lambda - \lambda^2) + \left( \frac{c}{l} + \frac{d}{l} \lambda \right) (a + b\lambda) = 0
\end{gather}

После раскрытия скобок и приведения подобных членов получаем характеристическое уравнение:

\begin{gather}
\lambda^4 - \left( d + \frac{b}{l} \right) \lambda^3 - \left( c + \frac{g}{l} + \frac{a}{l} \right) \lambda^2 + \frac{g}{l} d \lambda + \frac{g}{l} c = 0
\end{gather}

Запишем полином (19) в стандартной форме:

\begin{gather}
\lambda^4 + A_3 \lambda^3 + A_2 \lambda^2 + A_1 \lambda + A_0 = 0
\end{gather}

где

\begin{gather}
A_3 = -\left( d + \frac{b}{l} \right), \quad A_2 = -\left( c + \frac{g}{l} + \frac{a}{l} \right), \quad A_1 = \frac{g}{l} d, \quad A_0 = \frac{g}{l} c
\end{gather}

\begin{figure}[h!]
\centering
\includegraphics[width=0.7\textwidth]{ystal.jpg}
\caption{Я устал, я ухожу.}
\label{fig:risunok}
\end{figure}

куку








\end{document}
